\section{Forward-simulated short-premium scenarios}\label{sec:scenarios}

We forward-simulated 1{,}000 Goldman Sachs (GS) price paths over a 31-calendar-day horizon spanning 22 exchange sessions and observed that the model's $t=0$ Leisen-Reimer fair value sat essentially at the market mid on both legs (put: \$17.26 vs.\ \$16.51, edge $+\$0.75$; call: \$16.07 vs.\ \$16.09, edge $-\$0.02$; Supplementary Table~\ref{tab:gs_scenario}). The two contracts were real GS options pulled directly from the 2026-04-28 ladder capture: a short put at $K_p = \$890$ (market $\Delta = -0.295$, mid \$16.51, market IV $31.3\%$) and a short call at $K_c = \$970$ (market $\Delta = +0.328$, mid \$16.09, market IV $28.9\%$), both expiring 2026-05-29 with the underlying spot at $S_0 = \$926.38$. Each path drew from the JumpHMM marginal pretrained on GS returns. Conditional on the path, we evolved a per-contract Cox-Ingersoll-Ross variance $v_t$ with mean-reversion $\kappa = 15.0$, vol-of-vol $\sigma_v = 0.5$, and leverage coupling $\rho = -0.6$ to the per-contract target $\theta_i \cdot \psi_{\mathrm{NN}}((K/S_t)_{\mathrm{std}}, (T-t)_{\mathrm{std}})$, in which the per-ticker neural smile/term-structure surface entered through the drift target so the put and call each followed their own variance trajectory under a shared return shock. The two clocks stayed separate: the JumpHMM and the variance factor advanced one exchange session at a time on a $\Delta t = 1/252$ trading-year step, while the surface input and the lattice maturity used the calendar days remaining to expiry. At every trading day along every path we re-priced both contracts via 201-step Leisen-Reimer American with $\sigma_t^2(K, S_t) = v_t$, the contract's variance state from the SDE. We adopted Leisen-Reimer over the standard CRR lattice because the Peizer-Pratt construction places the strike $K$ at a fixed lattice position by design at every $(S, \sigma)$ configuration, which gave smooth Greeks via finite differences without the staircase aliasing that CRR's recombining tree produced when bumps moved $K$ across node boundaries. The seller received the market mid as the entry premium, and the model marked the position to its Leisen-Reimer fair value at every subsequent step; we assumed no scheduled earnings event during the window so the only sources of variation were diffusion, jumps, and the variance dynamics themselves. With the entry edge near zero, the simulated short was being sold at par with the model's own fair value at the entry capture and the realised path-conditional P\&L could be read directly as a return.

The forward-simulation engine is summarised in Algorithm~\ref{alg:forward_sim}. Each simulated exchange session along each path produced a $(S_t, \sigma_t)$ pair from which the LR pricer returned a fair value for the contract; the same $(S_t, \sigma_t)$ pair fed Supplementary Algorithm~\ref{alg:greeks} for the central finite-difference Greeks. The two algorithms shared the lattice constructor and the per-ticker $\psi_{\mathrm{NN}}$ surface, so the Greek bumps inherited the same shape function used for the base mark and did not require a separate calibration. For the scenarios reported in this section, we operated the variance process at the calibration-collapsed level $\bar\theta_i := \theta_i$ from \eqref{eq:theta_calibration}, with the regime-state and aggregate-mood multipliers in the full $\theta$ mechanism of \eqref{eq:theta} held inactive ($s_t \equiv s_0$, $\gamma = 0$); this isolates the contribution of the per-ticker $\psi_{\mathrm{NN}}$ surface to path-conditional pricing from the state-and-mood coupling that the framework also supports.

\begin{algorithm}[ht]
\caption{Forward simulation of path-conditional option premia}\label{alg:forward_sim}
\begin{algorithmic}[1]
\Require JumpHMM marginal for ticker $i$; per-ticker $\psi_{\mathrm{NN}}^{(i)}$ and calibration-collapsed level $\bar\theta_i := \theta_i$ from \eqref{eq:theta_calibration}; contract $(K, T)$; exchange sessions $\ell = 0,\ldots,L$ between the anchor date and expiry; calendar days to expiry $\tau_\ell$ at each session; CIR parameters $(\kappa, \sigma_v, \rho)$; spot $S_0$; risk-free rate $r$; number of paths $N_p$; LR tree depth $N_{\mathrm{LR}}$; trading-year step $\Delta t = 1/252$.
\Ensure $\{S_\ell^{(j)}, v_\ell^{(j)}, \sigma_\ell^{(j)}, P_\ell^{(j)}\}$ for $j = 1,\ldots,N_p$ and $\ell = 0, \ldots, L$.
\State \Comment{Operates at $\bar\theta_i$ from \eqref{eq:theta_calibration}; $s_t \equiv s_0$ and $\gamma = 0$ in this section}
\State Draw $N_p$ price paths $\{S_\ell^{(j)}\}_{\ell=0}^{L}$ from the JumpHMM marginal, one step per exchange session
\State $Z_{S,\ell}^{(j)} \gets$ the one-session log returns $\ln(S_\ell^{(j)}/S_{\ell-1}^{(j)})$ standardized over the whole ensemble
\For{$j = 1$ \textbf{to} $N_p$}
  \State $v_0^{(j)} \gets \bar\theta_i \cdot \psi_{\mathrm{NN}}^{(i)}\!\big(\ln(K/S_0),\, \ln \tau_0\big)$ \Comment{per-contract entry, eq.~\eqref{eq:v0}}
  \For{$\ell = 0$ \textbf{to} $L-1$}
    \State $\sigma_\ell^{(j)} \gets \sqrt{v_\ell^{(j)}}$ \Comment{per-contract IV, eq.~\eqref{eq:sigma_path}}
    \State $P_\ell^{(j)} \gets \mathrm{LR}_{N_{\mathrm{LR}}}\!\big(S_\ell^{(j)},\, K,\, \sigma_\ell^{(j)},\, \tau_\ell/365,\, r\big)$ \Comment{Leisen-Reimer American on the calendar clock}
    \State Draw $Z_\perp \sim \mathcal{N}(0,1)$ independent of the return shock
    \State $Z_{v,\ell}^{(j)} \gets \rho\, Z_{S,\ell+1}^{(j)} + \sqrt{1 - \rho^2}\, Z_\perp$ \Comment{return-variance coupling}
    \State $\theta_\ell^{(j)} \gets \bar\theta_i \cdot \psi_{\mathrm{NN}}^{(i)}\!\big(\ln(K/S_\ell^{(j)}),\, \ln \tau_\ell\big)$ \Comment{per-contract drift target}
    \State $v_{\ell+1}^{(j)} \gets \big| v_\ell^{(j)} + \kappa(\theta_\ell^{(j)} - v_\ell^{(j)})\,\Delta t + \sigma_v \sqrt{\max(v_\ell^{(j)},\,0)}\,\sqrt{\Delta t}\, Z_{v,\ell}^{(j)} \big|$
  \EndFor
  \State $P_L^{(j)} \gets \max(\phi (S_L^{(j)} - K),\, 0)$ \Comment{$\phi=+1$ call, $\phi=-1$ put}
\EndFor
\State \Return $\{S_\ell^{(j)}, v_\ell^{(j)}, \sigma_\ell^{(j)}, P_\ell^{(j)}\}$
\end{algorithmic}
\end{algorithm}

We rendered the share-price ensemble alongside the short-put and short-call premium ensembles (Fig.~\ref{fig:gs_paths}). On the worst-5\% paths (red, top row), GS drifted toward and through $K_p$, the short-put premium climbed many-fold above its entry price, and the position closed with a mean P\&L of $-\$90.73$ on those paths. On the top 1--5\% paths (green, bottom row; the wildest 1\% were excluded for legibility), GS pushed past $K_c$ and the short-call premium spiked despite the calendar pulling extrinsic value toward zero; mean P\&L on the top-5\% paths was $-\$125.62$, $1.38{\times}$ the mean loss on the worst put paths. On both legs the median path stayed inside the 25--75\% band and closed flat at the premium received; the short-premium trade made money on the typical path, with both legs carrying a long left tail.

\begin{figure}[ht]
\centering
\includegraphics[width=\textwidth]{sections/figures/gs/gs_short_paths.pdf}
\caption{Forward GS share-price paths and short option premium paths for the real GS 2026-05-29 contracts, 1{,}000 simulated trajectories over 31 calendar days and 22 trading transitions. Top row: share path with $K_p = \$890$ marked (left) and short-put premium (right), with the worst 1--5\% by terminal price overlaid in red. Bottom row: share path with $K_c = \$970$ marked (left) and short-call premium (right), with the top 1--5\% by terminal price overlaid in green; the most extreme 1\% are omitted from the overlay so the figure is not dominated by spike outliers, but the underlying summary statistics use the full top 5\%. Black solid: median; shaded band: 25--75\% IQR. Blue dashed lines on the option-price panels mark the market mids received at sale (\$16.51 put, \$16.09 call).}
\label{fig:gs_paths}
\end{figure}

The short-call leg lost asymmetrically because IV at the contract strike rose as $t \to T$ on every path (Fig.~\ref{fig:gs_iv}). The rise was not a model artefact: it was the empirical rise of the smile wings at short maturities that the per-ticker $\psi_{\mathrm{NN}}$ learned from the ladder corpus. As $(T-t)$ shrank at fixed $\log(K/S)$, the $\ln\tau$ input slid down a fitted term structure along which $\psi_{\mathrm{NN}}$ at off-the-money strikes grew monotonically toward expiry, and a Gaussian diffusion would have underpriced that wing without the inflation, since short-horizon returns carried much heavier tails than long-horizon returns, which were close to Gaussian by central-limit averaging. The variance $v_t$ tracked the rising drift target $\bar\theta_i \cdot \psi_{\mathrm{NN}}(\cdot)$ via mean-reversion, so the IV rise into expiry was inherited directly from the term-structure curvature of the per-ticker $\psi_{\mathrm{NN}}$ rather than from a separate level component. On the green tail $S_t$ pushed away from the call strike, pulling the standardised moneyness further into the wing and amplifying the rise; on the red tail the leverage coupling lifted $v_t$ at the same time as $S_t$ fell, so the two effects added.

\begin{figure}[ht]
\centering
\includegraphics[width=\textwidth]{sections/figures/gs/gs_iv_trajectories.pdf}
\caption{Path-conditional implied volatility at the put strike (left) and call strike (right) over the 31-calendar-day horizon for the GS scenario. Both panels showed IV rising into expiry: at fixed $\log(K/S)$, shrinking $(T-t)$ moved the contract down a term structure that the per-ticker network learned to be progressively higher and steeper in the wings at short horizons. Tail-binned paths (red, green) showed an additional rise driven by $S_t$ moving away from $K$ and, on the put side, leverage coupling lifting $v_t$ on down-paths.}
\label{fig:gs_iv}
\end{figure}

We decomposed the position risk into $\Delta$, $\Gamma$, and Vega along the same paths (Supplementary Fig.~\ref{fig:gs_greeks}), computed by central finite differences on the same 201-step Leisen-Reimer American pricer used for the base marks (spot bumps of $\pm 1.5\%$ for $\Delta$ and $\Gamma$; $\sigma$ bumps of $\pm 0.005$ for Vega-per-1\%-IV). The long-put $\Delta$ began near $-0.30$ and bifurcated: on the worst-5\% paths it ran to $-1$ as the put pinned ITM at expiry, while on the rest it ran to $0$; the long-call mirror image held at the call strike on the top tail. $\Gamma$ collapsed to zero on every path that drifted away from its strike but spiked sharply on paths that pinned to the strike near expiry, the standard reason traders avoid carrying short-premium positions through expiry. Vega decayed monotonically toward zero on every path (the $\sqrt{T-t}$ factor dominates), but the worst-5\% put paths and top-5\% call paths held materially more vega than the median through about $t = 20$, exposing the short-premium seller to extra IV-fluctuation risk on the same paths that were already losing on $\Delta$. The signs of the position-level Greeks for the short trader are the negatives of the long Greeks shown.

Both legs showed the characteristic short-premium asymmetry (Fig.~\ref{fig:gs_pnl}; moments in Supplementary Table~\ref{tab:gs_scenario}): most paths kept the full premium, and a long left tail held the loss-making paths. The short put kept its full \$16.51 premium on $73.7\%$ of paths and closed with a mean P\&L of $+\$4.65$; the short call kept its full \$16.09 premium on $68.8\%$ of paths and its mean P\&L of $-\$1.01$ was net negative, with the worst-case loss reaching $-\$277.29$, $1.57{\times}$ the worst put loss and $17{\times}$ the premium itself. The asymmetry between the two legs reflected the negative skew that GS's per-ticker $\psi_{\mathrm{NN}}$ encodes (higher IV at down-strikes than at equally-far up-strikes) together with the leverage coupling that lifted $v_t$ on down-paths; the put premium already reflected the down-tail risk and earned a fair return for taking it, while the call leg ran into a fatter empirical right tail than its symmetric Black--Scholes premium could cover.

\begin{figure}[ht]
\centering
\includegraphics[width=\textwidth]{sections/figures/gs/gs_short_pnl_distributions.pdf}
\caption{Terminal P\&L distributions per contract (1{,}000 paths, log-y) for the GS scenario. Both legs showed the same short-premium pattern: most paths kept the full premium, and a long left tail held the loss-making paths. The short call's left tail was substantially fatter than the put's despite a smaller premium received.}
\label{fig:gs_pnl}
\end{figure}

We compared the simulated premium-kept rates with the trader's delta-as-probability heuristic and found both legs within $3.2\%$ of the rule's prediction (Fig.~\ref{fig:gs_pnl}; Supplementary Table~\ref{tab:gs_scenario}). The heuristic approximates $\Pr[\mathrm{ITM\ at\ expiry}] \approx |\Delta|$ and so predicts premium-kept $\approx 1 - |\Delta|$ for a short held to expiry. The GS put's $|\Delta| = 0.295$ implied a $70.5\%$ rule-predicted retention against our simulated $73.7\%$, and the call's $|\Delta| = 0.328$ implied $67.2\%$ against our $68.8\%$ (deviations $+3.2\%$ and $+1.6\%$). The agreement was a non-trivial sanity check: the rule derives from the Black-Scholes identity $\Pr[\mathrm{ITM}] = N(d_2) \approx N(d_1) = |\Delta|$ at short horizons and near-zero drift, while our terminal probabilities emerged from a composition of the JumpHMM marginal, Heston variance dynamics, per-ticker $\psi_{\mathrm{NN}}$, and Leisen-Reimer American pricer, each calibrated to a different objective (return autocovariance, ladder RMSE, lattice convergence) and none calibrated against the rule. Both residuals were positive, which the scenario clock accounts for: the window spans 31 calendar days but only 22 exchange sessions, so the simulated paths accumulated the return variance of 22 trading steps while each contract's $|\Delta|$ was quoted against the full 31-day calendar maturity, and both strikes were therefore breached slightly less often than the heuristic implied.

\subsection{Cross-ticker validation: LLY}\label{sec:lly_scenario}

We re-ran the identical pipeline on a real LLY option pair from the same 2026-04-28 capture and observed two things: the qualitative short-premium asymmetry from GS reproduced cleanly on a ticker in a different sector at a higher IV regime, and the model's $t=0$ fair value sat \$1.35--\$1.80 below the market mid on both legs, a reversal of the near-par GS entry edges. The contract was a short put at $K_p = \$825$ (market $\Delta = -0.303$, mid \$23.30, market IV $44.4\%$) and a short call at $K_c = \$940$ (market $\Delta = +0.309$, mid \$20.76, market IV $44.0\%$), both expiring 2026-05-29 with LLY spot $S_0 = \$873.83$; every other framework choice was held identical to the GS scenario.

\begin{figure}[ht]
\centering
\includegraphics[width=\textwidth]{sections/figures/lly/lly_short_pnl_distributions.pdf}
\caption{LLY terminal P\&L distributions per contract (1{,}000 paths, log-y), parallel to Fig.~\ref{fig:gs_pnl}. Short put kept full premium on $81.2\%$ of paths (mean P\&L $+\$16.69$); short call kept full premium on $80.3\%$ of paths (mean P\&L $+\$11.42$) with a worst-case loss of $-\$212.60$, $1.64{\times}$ the worst put loss and $10.2{\times}$ the premium received. Full statistics in Supplementary Table~\ref{tab:lly_scenario}; share-price paths, IV trajectories, and Greek panels in the Supplementary Material.}
\label{fig:lly_pnl}
\end{figure}

The terminal P\&L distribution (Fig.~\ref{fig:lly_pnl}; Supplementary Table~\ref{tab:lly_scenario}) reproduced the GS pattern: short put kept full premium on $81.2\%$ of paths, short call on $80.3\%$, with the short-call worst-case loss again exceeding the worst put loss by $1.64{\times}$ despite the smaller premium received. The same asymmetric-tail pattern on a different sector and a different jump-state mixture showed that the negative skew came from the framework rather than from ticker-specific calibration. Both LLY legs ran above their delta-as-probability prediction by $\sim 11\%$, roughly four times the GS residual, and the same gap appeared in dollars as the $-\$1.35$/$-\$1.80$ entry-edge column of Supplementary Table~\ref{tab:lly_scenario}. The mechanism is an IV-regime mismatch: LLY had dropped $\sim 5\%$ on 04-28, so the contracts were captured against an elevated-IV state while the per-ticker $\psi_{\mathrm{NN}}$ encoded the time-averaged shape across fifteen capture dates. This particular failure mode---a static surface mis-pricing a single-day regime---motivates the low-rank daily adapter introduced below; we revisit the LLY 04-28 case study there and show that a rank-2 adapter closes the entry-edge gap to $-\$1.15$/$+\$0.46$ with 135 daily trainable parameters.

\subsection{Cross-sector validation across five additional tickers}\label{sec:expanded_scenarios}

We re-ran the identical pipeline on five additional tickers from the 31-ticker universe drawn to span ETF (SPY), Energy (XOM), Tech (AMD), low-liquidity Healthcare (PFE, the one ticker that fell below the 2{,}000-observation per-ticker fit threshold and therefore inherits the Healthcare sector $\psi_{\mathrm{NN}}$ surface), and large-cap Healthcare (JNJ), and observed that the premium-retention structure reproduced across all five names while the direction of the tail asymmetry did not (Supplementary Table~\ref{tab:expanded_scenarios}). Strikes for each ticker were pulled from the same 2026-04-28 capture at the closest $\sim 30\Delta$ $\sim$30-DTE pair, the Heston backbone $(\kappa, \sigma_v, \rho) = (15.0, 0.5, -0.6)$ was held identical to the GS and LLY scenarios, and each ticker drew from its own JumpHMM marginal with a ticker-prior CCGR anchor matching its sector long-run return ($8\%$/yr for ETF, Energy, Healthcare; $12\%$/yr for Tech). The premium-kept rates clustered between $62.5\%$ and $79.9\%$ across all five names, because the design 30-delta of each contract anchored each leg near the canonical $1 - |\Delta| \approx 70\%$ premium-retention target. The worst-case losses did not order consistently: the worst-call loss exceeded the worst-put loss on AMD ($1.45{\times}$: worst-call $-\$214.89$ vs worst-put $-\$148.68$) and on XOM ($1.09{\times}$), while SPY, PFE, and JNJ produced their larger loss on the put leg. Each worst-case figure is the single most extreme path out of $1{,}000$, so we read the retention rates rather than the tail ratios as the comparable cross-ticker statistic. PFE, the one ticker falling back to its sector $\psi_{\mathrm{NN}}$ rather than a per-ticker fit, sustained the smallest dollar tail losses simply because its premiums were small in absolute dollars, but it preserved the same percentage premium-kept structure as the per-ticker tickers, indicating that the sector-NN fallback is operationally viable for low-liquidity names. JNJ exhibited the LLY-style negative entry edge on the put leg (model fair value $-\$0.98$ below market mid), reflecting the same time-averaged-$\psi$ vs.\ capture-day-regime gap noted for LLY on 04-28; the same gap closed on JNJ's call leg ($+\$0.45$ above market) because the call sat at a different point on the $\psi$ surface. The pipeline therefore ran end to end across sectors, liquidity tiers, and volatility regimes spanning market IV from $\sim 17\%$ on SPY to $\sim 66\%$ on AMD, and it placed every leg near its delta-implied retention. Whether the heavier short-call tail seen on GS and LLY is a general property of the per-ticker smile geometry is not settled by five single-path extremes; separating that from sampling noise would require repeated seeds, which we did not run.

\subsection{Earnings-window scenario: INTC across the 04-22 print}\label{sec:earnings_window}

The cross-ticker scenarios above and the GS/LLY headline scenarios all sat in event-free windows by construction. We close this section by running the framework on its most common practical use case: a forward window that deliberately spans a scheduled earnings print. We anchored the scenario at INTC's 2026-04-13 capture ($S_0 = \$65.17$) and ran a forward simulation over 32 calendar days and 24 trading transitions to the contract's 2026-05-15 expiry, placing INTC's 2026-04-22 print nine calendar days into the window. The contract was a $30\Delta$ pair from the same capture: a short put at $K_p = \$60$ (mid \$3.28, market IV $78.2\%$) and a short call at $K_c = \$75$ (mid \$2.62, market IV $73.6\%$). We chose INTC because it carries the largest pre-print to post-print IV-regime gap in the fifteen-date corpus: its 04-22 earnings produced a $+2192\%$ EPS surprise driving a $\sim 30\%$ realised spot move from \$65.17 on 04-13 to \$84.54 on 04-28 (the same print isolated by the earnings-aware holdout as the residual signal no calendar-based indicator could absorb).

The simulator produced a terminal P\&L distribution dominated by middle-of-the-JumpHMM-marginal trajectories rather than by the realised post-print rally: short put kept full premium on $82.9\%$ of paths and short call on $90.9\%$, both materially above the $\sim 70\%$ retention typical of the same $30\Delta$ design on event-free windows and above every entry in Supplementary Table~\ref{tab:expanded_scenarios}. The model's $t=0$ Leisen-Reimer fair value sat $+\$0.22$ on the put and $+\$0.34$ on the call above market mids, indicating that the framework regarded both contracts as fairly priced from the perspective of its time-averaged surface. The empirical $\sim 30\%$ INTC rally would have placed the short call \$9.54 in-the-money at the realised path; the $90.9\%$ retention rate reflects the absence of that path in the JumpHMM marginal rather than a calibration success. This is the framework's clearest single-cell limitation under the static-$\psi$ operating point used throughout this section, and we return to it below in the adaptive recalibration section, where a low-rank daily adapter applied to INTC's 04-23 ladder slice reduces the slice RMSE from $68.6\%$ to $5.97\%$ with 135 trainable parameters.

